% =====================================================================
% Discussion and Conclusion -- revision 2026-10-10
% Claims checked against results/batches/noise_20260927_134251/summary,
% results/runs/20260920_005419_reyna_pre_surgery/supplement and
% results/batches/postcal_20261007_220027. "% CHANGED" marks a change of
% content; "% NEW" marks an added sentence or paragraph. Requires
% \addbibresource{refs_drug_simulation.bib} (now also holds vandeSande2026).
% =====================================================================
\FloatBarrier
\section{Discussion}
\label{sec:discussion}
% NEW: opening summary, as in other IJNMBE papers
We aimed to test whether a patient-specific LPM calibrated to a child's pre-closure catheterization can predict the same child's circulation after transcatheter VSD closure. The weight-scaled model, with seven parameters selected by global sensitivity analysis, reproduced the pre-closure record and predicted the post-closure systolic and mean arterial pressure within 7\%. The prediction did not reproduce the post-closure chamber volumes or the rise of diastolic pressure, and a second calibration to the post-closure record attributed most of the difference to a higher systemic resistance after closure. Varying the pre-closure record within its measurement uncertainty changed the calibrated parameters markedly and shifted the post-closure errors in a consistent direction. In the following sections, we discuss these errors under measurement uncertainty, compare the forward test with other post-intervention studies, and relate the sensitivity results and the closure response to clinical physiology.

\subsection{Pre- and post-closure errors under measurement uncertainty}
\label{sec:disc_noise}
Figure~\ref{fig:noise} compares the errors of the calibrated pre-closure model and of the post-closure prediction when the pre-closure record was varied within its measurement uncertainty. Before closure, the errors of the mean pressures and of $Q_p/Q_s$ stayed close to zero, with medians within 4\%, so calibration reproduced the record whichever noisy version it received. Only $Q_s$ was underestimated in every replicate (median $-16\%$), because calibration traded systemic flow against the pressures. After closure, the errors shifted in a consistent direction. Diastolic RFA pressure was underpredicted in every replicate (median $-17\%$), and mean PA pressure and RVEDV were overpredicted in 16 of 17 replicates (medians $+13\%$ and $+32\%$), whereas mean arterial pressure stayed centered on the measurement (median $0\%$). The reference prediction lay within the replicate spread for most quantities and at its upper edge for systolic PA pressure and RVEDV. Because these offsets persist whatever version of the record is used, they come from the model structure and from the parameters held fixed at closure rather than from measurement noise. The spread of the volume errors, from $-21\%$ to $+19\%$ for LVEDV, shows that a volume prediction from this record needs an uncertainty range.

% CHANGED: R_SVEN minimum is 0.35 of reference (about a third); bound counts from noise_parameters_long.csv
The calibrated parameters varied much more than the errors (Figure~\ref{fig:noise}A). PA compliance ranged from $-11\%$ to $+293\%$ of its reference value. Systemic arterial compliance and the discharge coefficient ranged from about half to almost twice their reference values, and systemic venous resistance from about a third to almost twice. The systemic resistance scale and RV passive elastance more than doubled in some replicates. The pulmonary resistance scale and RV passive elastance ended at or within 0.1\% of their lower bounds in six and four replicates. Records that differ only within measurement error are therefore reproduced equally well by markedly different parameter sets. Schiavazzi et al.\ reported the same behavior when estimating closed-loop parameters of single-ventricle circulations under measurement uncertainty~\cite{Schiavazzi2017}, and Kachabi et al.\ showed that pulmonary vascular parameters inferred from routine data carry wide uncertainty~\cite{Kachabi2025}. % NEW
Van de Sande et al.\ found the same practical non-identifiability in a closed-loop model fitted to orthostatic responses, where different parameter combinations reproduced similar blood pressure and heart rate traces across optimization runs~\cite{vandeSande2026}. Owashi et al.\ found relative errors of about 12\% when LV pressure was estimated from invasive data in twelve patients~\cite{Owashi2020}, which is of the same order as the pre-closure errors here. The present results add that the post-closure errors are not only larger but also shifted, which a single deterministic calibration cannot reveal. In line with recommendations for mechanistic cardiovascular models~\cite{Eck2016,Colebank2024Guidelines}, a prediction from a sparse pre-closure record should therefore be reported with the range produced by measurement uncertainty.

\begin{figure}[!htbp]
\centering
\includegraphics[width=\linewidth]{D3_noise_boxplots.pdf}
\caption{Measurement-noise simulation, 17 recalibrated replicates. (A) Each calibrated parameter divided by its value in the reference calibration (log scale). (B) Model minus measurement, as a percentage of the measurement, for the calibrated pre-closure model (blue) and the post-closure prediction (orange). Boxes show the median and interquartile range, points show the individual replicates, and diamonds show the reference calibration and its prediction. The post-closure $Q_p/Q_s$ equals 1 in every replicate because the defect is closed, so its error is $+11\%$ throughout and only the diamond is visible.}
\label{fig:noise}
\end{figure}

\subsection{Comparison with post-intervention prediction studies}
\label{sec:disc_prediction}
Few patient-specific lumped-parameter studies test a forward prediction against measurements taken after the intervention. In transcatheter aortic valve replacement, Khodaei et al.\ and Bahadormanesh et al.\ built separate models of the pre- and post-procedural states from the measurements of each state~\cite{Khodaei2021,Bahadormanesh2023}, so the post-procedural state was reproduced rather than predicted. Ooida et al.\ predicted the effect of mitral edge-to-edge repair from pre-procedural models but compared the result with the post-procedural state only qualitatively~\cite{Ooida2024}. Chen et al.\ predicted hemodynamics after virtual surgery in ten infants with borderline left ventricles and compared the predictions with the operation chosen clinically, not with measured post-operative values~\cite{Chen2026BLV}. Song et al.\ showed that post-operative conditions can differ markedly from the pre-operative state, which limits predictions that carry pre-operative conditions forward~\cite{Song2025}.

Against this background, the present study is a quantitative forward test. The seven post-closure catheter pressures, none of which entered calibration, were predicted with an RMSE of 6.9~mmHg, against 5.2~mmHg for the calibrated pre-closure pressures (Figure~\ref{fig:parity}), and systolic and mean arterial pressure were within 7\% of their measurements.
% NEW (keep only if the no-change sentence stays in Results 3.3)
Carrying the pre-closure measurements forward unchanged gave 6.2~mmHg, so for a shunt of this size the prediction adds little to the pressures themselves. Its value lies in the flows and volumes, which no pre-closure measurement provides.
The errors that remained match the limitation described by Song et al. The model keeps every vascular and chamber property at its pre-closure value, so it cannot follow the rise of diastolic pressure or the early LV unloading reported after transcatheter VSD closure~\cite{Yuan2025}. The post-closure calibration confirms this. With the same six parameters refitted, the pressure RMSE fell to 2.5~mmHg, mainly through a higher systemic resistance (Section~\ref{sec:results_post_calib}). Earlier VSD models simulated defects of different sizes without calibration to a patient~\cite{Colebank2026} or interpreted catheterization data without predicting a post-intervention state~\cite{Ferrero2026Shunts}, so this result has no direct counterpart in the VSD literature.

\begin{figure}[!htbp]
\centering
\includegraphics[width=0.55\linewidth]{M4_parity.pdf}
\caption{Model against measured catheter pressure. Filled circles, pre-closure calibration targets. Open circles, pre-closure pressures that were not targets. Triangles, post-closure prediction. Dashed line, identity.}
\label{fig:parity}
\end{figure}

\subsection{Sensitivity results against clinical physiology}
\label{sec:disc_gsa}
The parameters that dominated each output in the sensitivity screen (Figure~\ref{fig:gsa}) agree with the determinants that clinical physiology assigns to them. In a restrictive defect, the orifice limits flow and the pressure difference between the ventricles sets the shunt~\cite{Rao2024,Haas_2024}. Accordingly, $Q_p/Q_s$ depended on the discharge coefficient and on systemic resistance, which sets the LV pressure that drives the shunt, while pulmonary resistance had almost no effect ($S_T = 0.01$). Pulmonary resistance controlled the mean and diastolic PA pressure, and PA compliance controlled the systolic PA pressure. This is the separation of the steady and pulsatile parts of RV afterload described by Kempton et al.\ and by Gouwy et al.\ in children~\cite{Kempton2025,Gouwy2025}.
% CHANGED: PA compliance (0.14) ranks above systemic venous compliance (0.11) for mean RA pressure
Mean RA pressure depended mainly on RV passive elastance and, to a lesser extent, on PA and systemic venous compliance, in line with the view that RA pressure reflects RV filling and the mean systemic filling pressure~\cite{Persichini2022}. Two results agree only in part. Systolic arterial pressure depended more on resistance than on arterial compliance, whereas pulse-pressure physiology gives compliance a larger role~\cite{FloresGeronimo2021}. $Q_s$ depended on arterial resistance, partly because blood volume, its main physiological determinant, was not screened~\cite{Persichini2022}.

% CHANGED: claims softened (C_d stayed close to its clinical seed, so it is not an independent estimate)
These results have three practical implications. First, the calibrated discharge coefficient stayed close to its seed, which was computed from the measured LV--RV gradient, the shunt flow and the defect diameter. This suggests that the shunt of a restrictive defect can be characterized from routine measurements, since clinicians already estimate the interventricular gradient by Doppler~\cite{Uppu2022}. Second, because pulmonary resistance hardly affected $Q_p/Q_s$ in this restrictive defect, a change in pulmonary resistance would not be expected to change the shunt. In a nonrestrictive defect it would be expected to dominate, so the parameters to calibrate should be selected again for each defect type. Third, the active elastances controlled the end-systolic volumes and ejection fractions but none of the measured pressures and flows, as ventricular--arterial coupling predicts~\cite{Vriz2021}. A pre-closure catheterization therefore cannot inform contractility, and an echocardiographic volume measurement before closure is needed before the model can predict chamber volumes after closure. Lechuga et al.\ reached a similar conclusion in pulmonary hypertension, where the screen retained only the parameters that governed measured pressure and flow~\cite{Lechuga2026}, as recommended for mechanistic cardiovascular models~\cite{Colebank2024Guidelines}.

\subsection{Post-closure response against clinical physiology}
\label{sec:disc_closure}
Removing a restrictive left-to-right shunt is expected to lower pulmonary flow to the systemic flow, lower PA pressure, leave systemic and RA pressure unchanged or slightly higher, and unload the left heart~\cite{Rao2024,Yuan2025}. The predicted response was consistent with most of these changes (Figure~\ref{fig:closure}). Pulmonary flow fell by 13\% to equal systemic flow, and $Q_p/Q_s$ became 1.00. Mean, systolic and diastolic PA pressure fell, mean arterial pressure rose, and RA pressure stayed at about 5~mmHg, the same directions as in the measurements. The model changes were smaller than the measured ones for every pressure except RA pressure, as expected from a model whose vascular properties do not change at closure. For a shunt of this size ($Q_p/Q_s = 1.19$), small changes are expected, and both the simulated and the measured pressure changes were small.

\begin{figure}[!htbp]
\centering
\includegraphics[width=0.75\linewidth]{D1_closure_response.pdf}
\caption{Change at closure, as a percentage of the pre-closure value, in the model (bars) and in the measurements (diamonds). A missing diamond means that the quantity was not measured in both states.}
\label{fig:closure}
\end{figure}

Two predicted changes were not consistent with the expected physiology. First, the model predicted no unloading of the left heart. Mean LA pressure stayed at 7.9~mmHg and LVEDV rose from 39.8 to 41.1~mL, whereas LV dimensions in children normalize after transcatheter VSD closure, most of the change occurring in the first month~\cite{Yuan2025}. The cause is visible in the flows. Pulmonary venous return per minute fell by 13\%, but heart rate fell at the same time from 119 to 108~bpm, so the return per beat fell by only 4\% while diastole lengthened. With chamber properties held at their pre-closure values, the LV filled to a slightly larger volume. Second, the model predicted a small fall of diastolic RFA pressure where the measurement rose by 11~mmHg, a change in vascular tone that a model with fixed arterial properties cannot reproduce. This diastolic error ($-16$~mmHg) is nevertheless smaller than the disagreement between the three post-closure measurement sites, which ranged from 55~mmHg by cuff to 72~mmHg in the DAo.

% CHANGED: "sufficient explanation" -> "plausible explanation"; NEW: the corrected Fick value
Three measured changes were not consistent with the physiology of closure, and they explain the largest flow errors. Systemic flow fell from 3.42 to 2.54~L/min, although closure does not reduce systemic output, and $Q_p/Q_s$ fell to 0.90, which is not possible after complete closure without a right-to-left pathway. The pre-closure flows came from the Fick method and the post-closure flows from echocardiography. In pediatric catheterization, Fick output computed with an assumed oxygen consumption has a mean bias of $+27.5\%$ against thermodilution or MRI~\cite{Ralston2023}, and Fick and MRI flows agreed poorly in children after heart transplantation~\cite{Schramm2023}. Removing a bias of this size from the pre-closure value gives about 2.7~L/min, close to the post-closure 2.54~L/min. The change of method is therefore a plausible explanation for the measured fall in systemic flow and for the $+25\%$ error of the predicted $Q_s$. The measured fall of systolic RFA pressure, from 100 to 89~mmHg, also ran against the expected response and coincided with a rise of systolic DAo pressure from 91 to 104~mmHg, which points to the measurement site rather than to the circulation.

\subsection{Post-closure calibration as a basis for drug simulation}
\label{sec:disc_drug}
The two calibrations give one parameter set for each state of the same child. The difference between them, mainly a 65\% higher systemic resistance after closure, is the part of the closure response that the prediction could not represent (Table~\ref{tab:post_calib_params}), and the post-closure parameters describe the circulation in which post-operative drugs act. Milrinone is given to children after repair of congenital heart defects to prevent low cardiac output~\cite{Hoffman2003}, and its inotropic and vasodilating effects correspond to model parameters, namely active elastance and systemic and pulmonary resistance. Lumped-parameter and closed-loop models have already been used in this way. Casas et al.\ personalized a reduced-order model at rest and under dobutamine and recovered the expected rise in contractility~\cite{Casas2018}. Guyton-type models have generated virtual patients to study blood pressure regulation and its treatment~\cite{Moss2012}. CircAdapt digital twins have predicted the response to a device therapy~\cite{Koopsen2024}, and virtual surgery has been simulated in single-ventricle models~\cite{Corsini2014}. At the cellular level, in silico drug trials predicted clinical cardiotoxicity more accurately than animal models~\cite{Passini2017}, and electromechanical models have predicted inotropic drug effects~\cite{Margara2021}. Together, these studies describe the digital twin of precision cardiology~\cite{CorralAcero2020}.

Before such simulations can guide treatment, three steps are required. First, a drug must be linked to the parameters it changes through a pharmacokinetic--pharmacodynamic relation, and the simulated response must be compared with a measured one, as was done here for closure. Second, records that differ only within measurement error were reproduced by different parameter sets (Section~\ref{sec:results_noise}), and the compliances changed between states mainly through their seeds. A simulated drug response should therefore be reported with the range produced by parameter uncertainty~\cite{Schiavazzi2017,Colebank2024Guidelines}, and its credibility should be assessed for its context of use~\cite{Viceconti2021}. Third, a single patient cannot show how the response varies between children, which requires a cohort of calibrated models~\cite{Niederer2020,Niederer2021}.

\FloatBarrier
\subsection{Limitations}
% CHANGED: removed "No post-closure calibration was performed ..." (now done); NEW: input ranges, septal interaction, reflexes and anesthesia, post-closure compliance bounds
This study has several limitations. Only one patient was studied, so no claim of population-level performance can be made. The prediction describes the acute state only, because chamber and vascular properties are held at their pre-closure values, and it cannot represent the adaptation of the LV and of vascular tone over the hours and weeks after closure. The study protocol schedules measurements one day after closure and at the first outpatient visit, and calibrating time-dependent parameters to these data is the next step. The model has no mechanical interaction between the ventricles through the septum and no autonomic reflexes, and both states were measured under anesthesia, which may itself have changed vascular tone between the two measurements. No chamber volume was measured before closure, so the chamber parameters were not calibrated. A single pulmonary arterial compartment could not reproduce the narrow PA pulse pressure, which left systolic PA pressure $5.1\sigma$ above its measurement. Pre- and post-closure systemic flow came from different methods, Fick and echocardiography, and the three post-closure sites disagreed on diastolic pressure by 17~mmHg. In the post-closure calibration, the arterial compliances changed mainly through their clinical seeds. Finally, the sensitivity screen used 128 simulations stopped at 10 cycles, with input ranges that we set ourselves, and sampled each input independently, without the resistance--compliance coupling used in calibration. Its indices therefore rank the parameters but are not precise values.

\subsection{Future work}
\label{sec:disc_future}
The next step is to move from one child to a cohort. As more children complete the protocol of the prospective study, the same pipeline can be applied to each of them, so that the prediction error can be reported across patients rather than for a single case. A cohort would also allow defect types to be compared, since the parameters retained by the sensitivity screen are expected to differ between restrictive and nonrestrictive defects (Section~\ref{sec:disc_gsa}). With several patients, the calibration could move from independent deterministic fits to a Bayesian or hierarchical framework that shares information between children and reports each parameter with its posterior uncertainty~\cite{Schiavazzi2017,Colebank2024Guidelines}. Calibrated patients could then serve as the basis of a virtual cohort~\cite{Niederer2020}. Van de Sande et al.\ took the opposite route, fitting population-averaged responses first and naming individual fitting as their next step~\cite{vandeSande2026}. Both routes are needed before a model can describe the variability between children rather than one child.

These steps point towards a digital twin that supports, but does not replace, clinical decisions around VSD closure~\cite{CorralAcero2020,Niederer2021}. Such a twin would be updated as new data arrive: the pre-closure record before the procedure, the post-closure record after it, and the measurements scheduled one day after closure and at the first outpatient visit. Two calibrations of the same child are a first step in this direction, and the post-closure parameters could serve as the starting point for simulating post-procedural treatment (Section~\ref{sec:disc_drug}). Three requirements remain before clinical use. First, a calibration currently takes about 85~min, so a faster surrogate, for example a refined version of the polynomial chaos expansions built for the sensitivity screen, is needed for use during a clinical session. Second, every prediction should be reported with its uncertainty, as the measurement-noise simulation showed. Third, the credibility of the model must be established for a defined context of use, as required for in silico evidence~\cite{Viceconti2021}. Until these steps are taken, the model should be regarded as a research tool for understanding the closure response in an individual child.

\FloatBarrier
\section{Conclusion}
\label{sec:conclusion}
In this work, a closed-loop lumped-parameter model of pediatric VSD was scaled by body weight, screened by global sensitivity analysis, calibrated to one child's pre-closure catheterization and used without refitting to predict the same child's post-closure state. Calibration reduced the RMSE of the nine targets from 0.274 to 0.105. After closure, systolic and mean arterial pressure were predicted within 7\% and the seven catheter pressures with an RMSE of 6.9~mmHg. The largest errors were in systolic PA pressure, carried over from calibration, and in RV end-diastolic volume, which no pre-closure measurement constrains. Measurement uncertainty changed the calibrated parameters markedly and shifted the post-closure errors in a consistent direction. % NEW
A second calibration to the post-closure record reduced the pressure error to 2.5~mmHg and showed a higher systemic resistance after closure.

% CHANGED: "so its parameters can be read clinically" -> physiological interpretation
The sensitivity structure of the model matched the physiology of a restrictive VSD, so its parameters have a physiological interpretation. Adding a pre-closure echocardiographic volume measurement to calibration, measuring the patient again after closure, and studying more children are the next steps towards using such a model to plan VSD closure for the individual child.
